\section{Circulant matrices turn correlations into dot products}\label{ch14:sec:circulant}
Periodic correlations can be read off from a matrix. Write the sequence $a$ as the first row, and obtain each further row from the one above it by moving every entry one place to the right, the last entry wrapping around to the front. One more step after $n$ rows would bring back the first row, so we stop with an $n$-by-$n$ matrix, its rows and columns numbered from $0$ to $n-1$ (not from $1$ to $n$ as in Chapter~\ref{ch:linear}).

\par\noindent\begin{minipage}{\linewidth}
\begin{definition}[Circulant matrix]\label{ch14:def:circulant}
Let $a$ be a sign sequence of length $n$. The \emph{circulant matrix} of $a$ is the $n$-by-$n$ matrix $H$ with entries
\[
H_{ij}=a_{(j-i)\bmod n},\qquad0\le i,j\le n-1.
\]
A matrix of this form is called \emph{circulant}.
\end{definition}
\end{minipage}\par

Row $i$ is the sequence rotated $i$ places to the right. Such a rotation moves position $k$ to $(k+i)\bmod n$, so the entry in column $j$ started at position $j-i$, or $j-i+n$ if $j-i$ is negative; in both cases that is $(j-i)\bmod n$. For instance, in row $1$ the entry in column $0$ is $a_{(0-1)\bmod n}=a_{n-1}$, because $-1=(-1)\cdot n+(n-1)$ and remainders are never negative (Section~\ref{ch05:sec:euclid}). Row $0$ is the sequence itself, because $H_{0j}=a_j$, so every circulant matrix is the circulant matrix of its own first row. For $a=(1,1,1,-1)$,
\[
H=\begin{pmatrix}
1&1&1&-1\\
-1&1&1&1\\
1&-1&1&1\\
1&1&-1&1
\end{pmatrix}.
\]

As $j$ runs through $0,1,\ldots,n-1$, the index $(j-i)\bmod n$ also runs through $0,1,\ldots,n-1$, each value exactly once; with $n=4$ and $i=1$, the columns $j=0,1,2,3$ give the indices $3,0,1,2$. In the language of Section~\ref{ch03:sec:injsurj}, the map $j\mapsto(j-i)\bmod n$ is a bijection from $\{0,1,\ldots,n-1\}$ to itself, with inverse $k\mapsto(k+i)\bmod n$. Indeed, $(j-i)\bmod n+i$ is congruent to $j$ modulo $n$ (Section~\ref{ch06:sec:residues}), so it has the same remainder as $j$, which is $j$ itself because $0\le j\le n-1$; the other composition works in the same way. So each row of $H$ contains every entry of $a$ exactly once, and has squared length $a_0^2+\cdots+a_{n-1}^2=n$.

\subsection*{Dot products of rows are periodic correlations}
The dot product (Section~\ref{ch10:sec:inner}) of row $i$ with row $\ell$ is
\[
 \sum_{j=0}^{n-1}a_{(j-i)\bmod n}\,a_{(j-\ell)\bmod n}.
\]
For example, when $n=4$, row $0$ is $(a_0,a_1,a_2,a_3)$ and row $1$ is $(a_3,a_0,a_1,a_2)$, so their dot product is $a_0a_3+a_1a_0+a_2a_1+a_3a_2$, which is exactly $C_3$ as written out in Section~\ref{ch14:sec:periodic}. For the matrix $H$ above it is $-1+1+1-1=0$.

In general, put $k=(j-i)\bmod n$. By the bijection above, we may sum over $k$ instead of $j$: the terms are the same, in a different order. For the second factor, $k\equiv j-i\pmod n$ gives
\[
 j-\ell=(j-i)+(i-\ell)\equiv k+(i-\ell)\pmod n.
\]
Congruent integers have the same remainder, so $(j-\ell)\bmod n=(k+i-\ell)\bmod n$. Now let $s=(i-\ell)\bmod n$. Then $s\equiv i-\ell\pmod n$, so $(k+i-\ell)\bmod n=(k+s)\bmod n$ for the same reason, and the dot product becomes
\[
 \sum_{k=0}^{n-1}a_k\,a_{(k+s)\bmod n}=C_s.
\]

\par\noindent\begin{minipage}{\linewidth}
\begin{proposition}\label{ch14:prop:rows}
Let $a$ be a sign sequence of length $n$, and let $H$ be its circulant matrix.
\begin{enumerate}
\item[(a)] For all rows $i$ and $\ell$, the dot product of row $i$ with row $\ell$ equals $C_{(i-\ell)\bmod n}(a)$.
\item[(b)] Any two different rows of $H$ are orthogonal exactly when $C_s(a)=0$ for every shift $s$ with $1\le s\le n-1$.
\end{enumerate}
\end{proposition}
\end{minipage}\par

\begin{proof}
Part (a) is the calculation above. For part (b), note first that if $i\ne\ell$, then $i-\ell$ lies between $-(n-1)$ and $n-1$ and is not zero, so it is not a multiple of $n$, and the shift $(i-\ell)\bmod n$ is not zero. So by part (a), the dot product of two different rows is some $C_s$ with $1\le s\le n-1$, and if all these $C_s$ are zero, different rows are orthogonal. Conversely, if any two different rows are orthogonal, then each $C_s$ with $1\le s\le n-1$ is zero, because by part (a) it is the dot product of the different rows $i=s$ and $\ell=0$, as $(s-0)\bmod n=s$.
\end{proof}

\subsection*{Circulant Hadamard matrices}
Recall from Section~\ref{ch10:sec:product} that a square matrix $H$ with entries $\pm1$ is a \emph{Hadamard matrix} when $HH^{\mathsf T}=nI$, where $n$ is the number of rows; this number $n$ is called the \emph{order} of the matrix. The diagonal entries of this identity say that every row has squared length $n$, which is automatic for entries $\pm1$, and the others say that different rows are orthogonal. So a matrix with entries $\pm1$ is a Hadamard matrix exactly when any two different rows are orthogonal.

A \emph{circulant Hadamard matrix} is a Hadamard matrix that is also circulant. Since every circulant matrix is the circulant matrix of its own first row, Proposition~\ref{ch14:prop:rows}(b) and the previous paragraph give:
\begin{quote}
A circulant Hadamard matrix of order $n$ exists exactly when some sign sequence of length $n$ has $C_s=0$ for every shift $s$ with $1\le s\le n-1$.
\end{quote}
Since $(1,1,1,-1)$ has $C_1=C_2=C_3=0$, the matrix $H$ above is a circulant Hadamard matrix of order four.

Habit~6 (Section~\ref{ch06:sec:order}) asks what a change of description preserves. Here nothing is lost: we have the same question in two languages, one suggesting sums and shifts, the other dot products and orthogonality. The next two sections use each in turn.
